\chapter{Metodologia sperimentale}
\label{cap:metodologia}

\section{Ambiente sperimentale}
\label{sec:ambiente-sperimentale}

Le campagne sperimentali sono state eseguite sul server di calcolo del
Dipartimento. La configurazione hardware e software rilevante ai fini delle
misure è riassunta in \cref{tab:ambiente-sperimentale}.

\begin{table}[H]
\centering
\footnotesize
\caption{Configurazione dell'ambiente sperimentale.}
\label{tab:ambiente-sperimentale}
\begin{tabular}{@{}ll@{}}
\toprule
\textbf{Componente} & \textbf{Configurazione} \\
\midrule
CPU & \TODO{modello CPU} \\
Socket / core & \TODO{numero socket e core} \\
Memoria principale & \TODO{quantità di RAM} \\
\midrule
GPU & NVIDIA Quadro RTX 5000 \\
Memoria GPU & \SI{15.56}{\gibi\byte} \\
Compute Capability & 7.5 (Turing) \\
Streaming Multiprocessor (SM) & 48 \\
\midrule
Compilatore GNU & GCC 13.3.0 \\
Open MPI & 5.0.9 \\
CUDA Toolkit & 12.8 \\
Architettura CUDA & \code{sm\_75} \\
\bottomrule
\end{tabular}
\end{table}

L'ambiente software viene predisposto tramite il sistema di moduli del server,
utilizzando versioni fissate della toolchain GNU, di Open MPI e di CUDA:

\begin{verbatim}
module purge
module load gnu/13.3.0
module load openmpi/5.0.9
module load cuda/12.8
\end{verbatim}

In questo modo tutte le campagne vengono compilate ed eseguite utilizzando la
stessa toolchain software.

Le implementazioni CPU sono compilate con ottimizzazione \code{-O3} e
ottimizzazioni specifiche per l'architettura host tramite
\code{-march=native}; viene inoltre abilitata la contrazione delle operazioni
floating-point mediante \code{-ffp-contract=fast}. Le implementazioni CUDA sono
compilate con \code{-O3} per l'architettura \code{sm\_75}, corrispondente alla
Compute Capability 7.5 della GPU utilizzata. La precisione numerica, singola o
doppia, viene selezionata a tempo di compilazione.

Per le esecuzioni MPI multi-processo i rank vengono vincolati ai core mediante
le opzioni \code{--bind-to core --map-by core}. Il placement esplicito riduce
la variabilità dovuta alla migrazione dei processi e rende più controllabile la
località degli accessi alla memoria durante il confronto tra configurazioni.

Il nodo dispone di una singola GPU e tutti i rank CUDA utilizzano il device
\code{0}. Di conseguenza, quando più rank MPI eseguono contemporaneamente un
backend CUDA, essi condividono la stessa GPU e i tempi del kernel possono
includere effetti di contesa e alternanza tra i diversi contesti CUDA.
Per questo motivo le campagne finalizzate al confronto diretto tra kernel GPU
e alla caratterizzazione della loro efficienza computazionale vengono eseguite
con un solo rank MPI ($P=1$). Le eventuali esecuzioni CUDA con $P>1$ descrivono
invece il comportamento del sistema in presenza di condivisione della singola
GPU e non uno strong scaling su più acceleratori.

Per ogni campagna il runner registra inoltre nei metadati il comando eseguito,
i principali parametri sperimentali, il commit Git, lo stato del working tree,
le versioni di CUDA e Open MPI e le informazioni sulla GPU. Tali informazioni
consentono di associare ogni insieme di misure alla configurazione software con
cui è stato prodotto.

\section{Protocollo di misura}
\label{sec:protocollo-misura}

Le misure vengono raccolte durante l'esecuzione del prodotto distribuito,
separando le tre fasi principali dell'algoritmo: distribuzione del blocco
locale di $X$, calcolo locale e riduzione dei contributi di $Y$.

Prima di ogni repetition misurata viene eseguita una barriera sul comunicatore
della griglia MPI, in modo che tutti i rank inizino l'operazione dalla stessa
fase dell'esecuzione. La barriera precede l'avvio del cronometro e non è quindi
inclusa nei tempi riportati.

Per ciascun rank, una repetition segue la sequenza

\begin{equation}
\text{\code{MPI\_Bcast}}
\;\longrightarrow\;
\text{\code{local\_gemm}}
\;\longrightarrow\;
\text{\code{MPI\_Reduce}} .
\end{equation}

Il broadcast rende disponibile a tutti i processi interessati il blocco di
$X$ necessario al calcolo locale; il risultato parziale prodotto da
\code{local\_gemm} viene successivamente combinato mediante la reduce.
Le due operazioni collettive utilizzano rispettivamente i comunicatori di
colonna e di riga della griglia bidimensionale.

I quattro istanti temporali vengono acquisiti mediante \code{MPI\_Wtime}:

\begin{align}
t_0 &:\quad \text{inizio del broadcast}, \\
t_1 &:\quad \text{fine del broadcast}, \\
t_2 &:\quad \text{fine della fase locale}, \\
t_3 &:\quad \text{fine della reduce}.
\end{align}

Da essi vengono ricavati i tempi

\begin{align}
t_{\mathrm{bcast}} &= t_1-t_0, \\
t_{\mathrm{local}} &= t_2-t_1, \\
t_{\mathrm{reduce}} &= t_3-t_2, \\
t_{\mathrm{total}} &= t_3-t_0.
\end{align}

Il tempo $t_{\mathrm{total}}$ rappresenta quindi il tempo wall-clock
dell'intera invocazione distribuita, dall'inizio del broadcast alla conclusione
della reduce. Le operazioni di preparazione del problema e di inizializzazione
del backend vengono eseguite prima delle repetition e non appartengono a tale
intervallo temporale.
\paragraph{Misura della fase CUDA.}
Nei backend CPU la fase \code{local\_gemm} coincide con il calcolo locale.
Nei backend CUDA, invece, la stessa fase comprende il trasferimento del blocco
di $X$ dalla memoria host alla GPU, l'esecuzione del kernel e il trasferimento
del risultato locale $Y$ dalla GPU alla memoria host:

\begin{equation}
t_{\mathrm{local}}
\simeq
t_{\mathrm{H2D},X}
+
t_{\mathrm{kernel}}
+
t_{\mathrm{D2H},Y}
+
t_{\mathrm{overhead}} .
\end{equation}

I tre intervalli GPU principali vengono misurati separatamente mediante
CUDA events. In particolare vengono registrati il tempo del trasferimento
H2D di $X$, il tempo di esecuzione del kernel e il tempo del trasferimento
D2H di $Y$. La sincronizzazione associata al completamento dell'ultimo
trasferimento garantisce che la fase locale sia conclusa prima che venga
eseguita la successiva \code{MPI\_Reduce}.

La copia iniziale della matrice locale $A$ verso la memoria della GPU avviene
invece durante l'inizializzazione del backend e non viene ripetuta a ogni
invocazione del prodotto.


\section{Definizione della metrica di prestazione}
\label{sec:metrica-prestazione}

I tempi descritti nella sezione precedente vengono misurati separatamente da
ciascun rank MPI. Poiché l'esecuzione distribuita può proseguire solo alla
velocità del processo più lento, il tempo associato a una repetition viene
definito considerando il massimo tra i rank.

Indicando con $i$ la repetition e con $p$ il rank MPI, il tempo ufficiale
locale della repetition è definito in modo diverso per i backend CPU e CUDA.

Per i backend CPU:

\begin{equation}
T^{\mathrm{CPU}}_{i,p}
=
t_{\mathrm{bcast},i,p}
+
t_{\mathrm{local},i,p}
+
t_{\mathrm{reduce},i,p}.
\label{eq:t-official-cpu-local}
\end{equation}

Poiché in questo caso l'intera fase locale coincide con il calcolo eseguito
dalla CPU, tale quantità corrisponde anche al tempo wall-clock
$t_{\mathrm{total}}$ misurato dall'inizio del broadcast alla fine della reduce.

Per i backend CUDA la metrica richiesta esclude invece i trasferimenti tra
memoria host e device. La fase locale viene pertanto rappresentata dal solo
tempo di esecuzione del kernel:

\begin{equation}
T^{\mathrm{CUDA}}_{i,p}
=
t_{\mathrm{bcast},i,p}
+
t_{\mathrm{kernel},i,p}
+
t_{\mathrm{reduce},i,p}.
\label{eq:t-official-cuda-local}
\end{equation}

Il tempo rappresentativo della repetition $i$ è quindi

\begin{equation}
T_i = \max_{0 \leq p < P} T_{i,p},
\label{eq:t-official-repetition}
\end{equation}

dove $P$ è il numero complessivo di processi MPI.

La composizione delle diverse fasi viene effettuata localmente prima
dell'operazione di massimo. In particolare, nel caso CUDA viene calcolato

\begin{equation}
\max_p
\left(
t_{\mathrm{bcast},p}
+
t_{\mathrm{kernel},p}
+
t_{\mathrm{reduce},p}
\right),
\end{equation}

e non la somma

\begin{equation}
\max_p t_{\mathrm{bcast},p}
+
\max_p t_{\mathrm{kernel},p}
+
\max_p t_{\mathrm{reduce},p}.
\end{equation}

La seconda espressione potrebbe infatti combinare tempi appartenenti a rank
differenti e non rappresenterebbe il percorso temporale di alcun processo
reale.

Dopo aver ottenuto un valore $T_i$ per ciascuna delle $R$ repetition misurate,
il tempo medio utilizzato per valutare le prestazioni è

\begin{equation}
\overline{T}
=
\frac{1}{R}
\sum_{i=1}^{R} T_i .
\label{eq:t-official-mean}
\end{equation}

Il prodotto tra una matrice $A \in \mathbb{R}^{M \times N}$ e un
multivettore $X \in \mathbb{R}^{N \times k}$ richiede approssimativamente

\begin{equation}
2MNk
\end{equation}

operazioni floating-point. La prestazione viene quindi espressa come

\begin{equation}
\mathrm{GFLOP/s}
=
\frac{2MNk}{\overline{T}} \cdot 10^{-9}.
\label{eq:gflops}
\end{equation}

Questa quantità costituisce la metrica principale utilizzata per il confronto
tra le diverse configurazioni sperimentali.

\paragraph{Metriche ausiliarie.}
Oltre alla metrica complessiva, il programma calcola una misura della sola fase
computazionale. Per i backend CPU essa utilizza $t_{\mathrm{local}}$, mentre
per i backend CUDA utilizza $t_{\mathrm{kernel}}$:

\begin{equation}
T_{\mathrm{compute}} =
\begin{cases}
t_{\mathrm{local}}, & \text{CPU}, \\[2pt]
t_{\mathrm{kernel}}, & \text{CUDA}.
\end{cases}
\end{equation}

La corrispondente prestazione è riportata come \code{gflops\_compute}. Nei
backend CUDA la stessa quantità è riportata anche come
\code{gflops\_kernel}.

La colonna \code{gflops} rimane invece la metrica di riferimento, poiché
comprende anche il costo delle comunicazioni MPI previste dal prodotto
distribuito.



\section{Costi esclusi dalla metrica ufficiale}
\label{sec:costi-esclusi}

La metrica di prestazione definita nella sezione precedente considera
esclusivamente il prodotto distribuito vero e proprio. Restano quindi esclusi
dal tempo ufficiale i costi di preparazione del problema, l'I/O e, nel caso dei
backend CUDA, le operazioni di trasferimento tra memoria host e memoria device.

In particolare, il preprocessing comprende l'allocazione dei buffer locali,
la generazione o distribuzione iniziale dei dati, l'eventuale conversione del
layout di $X$ e l'inizializzazione del backend. Tali operazioni vengono
eseguite prima delle repetition misurate e non contribuiscono quindi a
$T_{\mathrm{official}}$.

Per i backend CUDA, anche il trasferimento iniziale della matrice locale $A$
verso la memoria della GPU appartiene alla fase di setup e viene eseguito una
sola volta. Durante ciascuna invocation di \code{local\_gemm} vengono invece
eseguiti il trasferimento H2D del blocco di $X$, il kernel CUDA e il
trasferimento D2H del risultato locale $Y$.

La metrica ufficiale utilizza, per la fase locale CUDA, il solo tempo del
kernel. Di conseguenza i trasferimenti H2D e D2H, insieme all'eventuale
overhead residuo lato host/runtime, non entrano in $T_{\mathrm{official}}$.
Essi vengono tuttavia misurati separatamente, così da poter distinguere il
costo computazionale del kernel dal costo complessivo necessario a utilizzare
l'acceleratore.

Il programma registra inoltre il tempo wall-clock completo della fase locale
CUDA:

\begin{equation}
t_{\mathrm{local}}
\simeq
t_{\mathrm{H2D},X}
+
t_{\mathrm{kernel}}
+
t_{\mathrm{D2H},Y}
+
t_{\mathrm{overhead}} ,
\end{equation}

dove $t_{\mathrm{overhead}}$ rappresenta un termine residuo ottenuto per
differenza e comprende costi host/runtime non attribuiti direttamente ai tre
intervalli misurati mediante CUDA events. Tale quantità non deve quindi essere
interpretata come una misura diretta del solo costo di kernel launch.

Anche l'I/O necessario alla produzione dei file di output viene mantenuto fuori
dai tempi del prodotto. In questo modo le scritture su file non influenzano le
misure di prestazione, pur rimanendo disponibili per la successiva analisi dei
risultati.

La separazione tra tempo ufficiale e costi accessori consente quindi di
utilizzare una metrica coerente per il confronto tra implementazioni, senza
rinunciare alla possibilità di valutare il costo reale del trasferimento dei
dati e dell'inizializzazione dell'acceleratore.

\section{Ripetizioni e stabilità statistica}
\label{sec:stabilita-misure}

Ogni configurazione sperimentale viene eseguita più volte all'interno della
stessa esecuzione del programma. Prima delle repetition utilizzate per le
statistiche vengono effettuate alcune invocazioni di warm-up, i cui tempi non
vengono registrati. Il numero di warm-up e di repetition è definito
esplicitamente nella configurazione della singola campagna.

Le invocazioni di warm-up hanno lo scopo di portare l'esecuzione in una
condizione più rappresentativa prima della raccolta dei campioni, evitando che
la prima misura sia influenzata da costi transitori. Al termine dei warm-up
vengono eseguite le $R$ repetition misurate secondo il protocollo descritto
nelle sezioni precedenti.

Prima di ciascuna repetition viene eseguita una \code{MPI\_Barrier} sul
comunicatore della griglia. La barriera è esterna all'intervallo cronometrato e
serve esclusivamente ad allineare l'inizio della nuova repetition tra i rank.


Per ciascuna repetition viene ottenuto il tempo globale $T_i$ definito in
\cref{eq:t-official-repetition}. L'insieme

\begin{equation}
\{T_1,T_2,\ldots,T_R\}
\end{equation}

costituisce il campione utilizzato per l'analisi statistica della
configurazione.

Il valore principale riportato nei confronti prestazionali è la media
aritmetica

\begin{equation}
\overline{T}
=
\frac{1}{R}
\sum_{i=1}^{R} T_i .
\end{equation}

Alla media vengono affiancate misure di dispersione e statistiche descrittive,
tra cui deviazione standard campionaria, mediana e minimo. La deviazione
standard è calcolata come

\begin{equation}
s =
\sqrt{
\frac{1}{R-1}
\sum_{i=1}^{R}
(T_i-\overline{T})^2
}.
\end{equation}

Per confrontare la variabilità di misure aventi scale temporali differenti
viene inoltre utilizzato il coefficiente di variazione

\begin{equation}
CV =
100 \frac{s}{\overline{T}},
\label{eq:cv}
\end{equation}

espresso in percentuale.

Il numero di repetition non viene considerato, da solo, una garanzia di
stabilità della misura. La dispersione viene invece verificata a posteriori
attraverso il coefficiente di variazione e, nei casi in cui questo evidenzi
una variabilità significativa, mediante l'ispezione dei singoli campioni.
Differenze prestazionali dello stesso ordine della dispersione osservata non
vengono considerate sufficienti, da sole, per trarre conclusioni sul confronto
tra due configurazioni.

Oltre ai valori aggregati, il programma può produrre un file CSV contenente
una riga per ciascuna repetition misurata. I campioni raw consentono di
individuare eventuali outlier, variazioni sistematiche nel corso
dell'esecuzione o altri comportamenti che potrebbero essere nascosti dalla
sola media.

Le invocazioni di warm-up non vengono incluse nei file raw né nelle
statistiche finali.

Nelle campagne principali vengono utilizzate tipicamente 20 repetition
precedute da 5 warm-up; per le campagne con un numero maggiore di
configurazioni vengono impiegati, dove previsto, 10 campioni dopo 3 warm-up.
I valori effettivi sono comunque parte della configurazione della singola
campagna e vengono registrati nei relativi metadati.

\section{Disegno delle campagne sperimentali}
\label{sec:campagne-sperimentali}

Le campagne sperimentali sono organizzate in modo da isolare, per quanto
possibile, l'effetto delle principali variabili del problema: dimensione e
forma della matrice, numero di colonne del multivettore, implementazione del
kernel locale, numero di processi MPI, forma della griglia bidimensionale e
precisione numerica.

Ogni campagna modifica quindi un insieme limitato di parametri mantenendo
costanti gli altri. Le configurazioni utilizzate sono definite tramite file
\code{.conf}, in modo da rendere espliciti e riproducibili i parametri di ogni
esperimento.


\subsection{C1 --- Dimensione del problema e rapporto $M{:}N$}
\label{sec:campagna-c1}

La prima campagna studia l'effetto della dimensione globale del problema e
della forma della matrice. Vengono considerati tre rapporti tra le dimensioni:

\begin{equation}
M:N \in \{1:1,\;3:1,\;1:2\},
\end{equation}

includendo quindi sia matrici quadrate sia configurazioni rettangolari.

Per ciascun rapporto vengono utilizzate quattro taglie crescenti. I valori di
$M$ e $N$ sono scelti in modo che, a parità di livello, il prodotto $MN$ sia
approssimativamente equivalente tra i tre rapporti. Ciò permette di confrontare
problemi con una quantità di lavoro simile ma con una diversa geometria della
matrice.

La campagna viene eseguita sia per il backend CPU sia per un backend CUDA,
utilizzando un singolo processo MPI ($P=1$), in modo da isolare gli effetti
della dimensione e della forma del problema dalla comunicazione MPI.
Vengono considerati $k=8$ e $k=32$, rappresentativi di due differenti
dimensioni del multivettore.


\subsection{C2 --- Sweep su $k$ e confronto dei backend}
\label{sec:campagna-c2}

La seconda campagna confronta le diverse implementazioni mantenendo fisse le
dimensioni del problema e variando il numero di colonne $k$ del multivettore.

Sono innanzitutto considerati i valori richiesti dalla traccia,
$k \in \{3,6,8,20,32\}$. A questi vengono aggiunti alcuni valori non
specializzati, utilizzati per verificare il comportamento del percorso
generico dell'implementazione.

Il problema viene mantenuto fisso a $M=N=10000$ e viene utilizzato un singolo
processo MPI ($P=1$). In questo modo il confronto riflette principalmente le
differenze tra le implementazioni locali, senza introdurre il costo delle
comunicazioni tra processi.

Le implementazioni confrontate comprendono il backend CPU
\kernelname{scheme\_a}, i kernel CUDA
\kernelname{cuda\_naive},
\kernelname{cuda\_warp},
\kernelname{cuda\_warp\_smem} e
\kernelname{cuda\_warp\_tiled}, oltre all'implementazione basata su cuBLAS,
utilizzata come riferimento.

\label{sec:campagna-c3}

La campagna C3 è progettata come confronto controllato tra la versione
specializzata e la versione generica del kernel CPU.

Le due configurazioni utilizzano gli stessi valori di $M$, $N$, $k$ e lo stesso
numero di processi; l'unica differenza rilevante è l'abilitazione o meno della
specializzazione a compile-time per $k$.

Questo consente di attribuire le differenze osservate direttamente alla
specializzazione, evitando il confronto tra problemi con una quantità di
lavoro differente.

La campagna viene eseguita con $M=N=10000$ e $P=1$.

\subsection{C4 --- Strong scaling}
\label{sec:campagna-c4}

La campagna C4 studia lo strong scaling dell'implementazione MPI.

La dimensione globale del problema viene mantenuta costante a

\begin{equation}
M=N=20000,
\end{equation}

mentre il numero di processi varia secondo

\begin{equation}
P \in \{1,2,4,8,16,20\}.
\end{equation}

Per ogni valore di $P$ viene utilizzata una griglia bidimensionale prossima
alla forma quadrata, in modo da evitare di introdurre intenzionalmente
configurazioni fortemente sbilanciate.

Vengono considerati $k=3$ e $k=32$, così da osservare il comportamento dello
scaling per due differenti dimensioni del multivettore.

Lo speedup di strong scaling viene calcolato rispetto alla configurazione a un
processo della stessa campagna:

\begin{equation}
S(P)=\frac{T(1)}{T(P)},
\end{equation}

mentre l'efficienza parallela è

\begin{equation}
E(P)=\frac{S(P)}{P}.
\end{equation}

\subsection{C5 --- Weak scaling}
\label{sec:campagna-c5}

La campagna C5 studia il weak scaling mantenendo approssimativamente costante
la quantità di dati assegnata a ciascun processo.

Per una griglia $P_r \times P_c$ vengono scelte le dimensioni globali secondo

\begin{equation}
M = 5000 P_r,
\qquad
N = 5000 P_c.
\end{equation}

Di conseguenza ciascun processo opera, salvo gli effetti della decomposizione,
su un blocco locale di circa

\begin{equation}
5000 \times 5000.
\end{equation}

All'aumentare del numero di processi cresce quindi la dimensione globale del
problema, mentre il lavoro computazionale locale nominale rimane
approssimativamente costante.

L'andamento del tempo totale e la sua scomposizione nelle componenti di
broadcast, calcolo locale e reduce permettono di osservare come il costo delle
comunicazioni e gli effetti del sistema evolvano all'aumentare di $P$.

\subsection{C6 --- Forma della griglia MPI}
\label{sec:campagna-c6}

La campagna C6 isola l'effetto della forma della griglia bidimensionale.

Per un numero complessivo di processi fissato vengono considerate tutte le
fattorizzazioni ordinate

\begin{equation}
P=P_rP_c.
\end{equation}

Ad esempio, per $P=8$ vengono confrontate le configurazioni

\begin{equation}
1\times8,\quad
2\times4,\quad
4\times2,\quad
8\times1.
\end{equation}

La dimensione globale del problema e il numero complessivo di processi
rimangono invariati all'interno di ciascun confronto; cambia solamente la
ripartizione tra righe e colonne della griglia MPI.

La campagna consente di valutare sperimentalmente l'effetto della
decomposizione bidimensionale sui costi di broadcast di $X$ e reduce di $Y$,
confrontando anche le configurazioni degeneri $1\times P$ e $P\times1$ con
quelle maggiormente bilanciate.


\subsection{C7 --- Precisione numerica}
\label{sec:campagna-c7}

L'ultima campagna studia l'effetto della precisione numerica sui backend CUDA.
Il confronto viene effettuato tra FP64 e FP32 mantenendo invariati, per ogni
coppia di misure, dimensione del problema, valore di $k$, kernel e
configurazione di lancio.

La campagna è suddivisa in due parti complementari.

\paragraph{C7.1 --- Prestazioni FP32 e FP64.}

I principali backend CUDA vengono compilati ed eseguiti sia in singola sia in
doppia precisione, con $M=N=10000$ e $P=1$.

Il confronto viene effettuato utilizzando la stessa metrica temporale per le
due precisioni. In particolare, quando l'obiettivo è analizzare il solo
comportamento del kernel CUDA, vengono confrontati i rispettivi valori di
$t_{\mathrm{kernel}}$.

Lo speedup dovuto alla precisione può quindi essere espresso come

\begin{equation}
S_{32/64}
=
\frac{t_{\mathrm{kernel}}^{FP64}}
     {t_{\mathrm{kernel}}^{FP32}}.
\end{equation}

\paragraph{C7.2 --- Utilizzo delle risorse CUDA.}

La seconda parte non costituisce un benchmark temporale, ma una
caratterizzazione delle risorse richieste dai kernel.

Le compilazioni vengono effettuate utilizzando l'output di \code{ptxas}, dal
quale vengono ricavati il numero di registri per thread, eventuali spill,
stack frame e shared memory statica.

Questi dati vengono utilizzati per correlare le differenze osservate tra FP32
e FP64 con la pressione sulle risorse hardware, senza considerarli una misura
diretta dell'occupancy effettivamente raggiunta durante l'esecuzione.

Nel complesso, le campagne sono organizzate in modo da separare gli effetti
delle diverse variabili: C1--C3 caratterizzano principalmente il comportamento
computazionale delle implementazioni locali, C4--C6 studiano il comportamento
della decomposizione MPI, mentre C7 approfondisce l'effetto della precisione
numerica sui backend CUDA.

Le conclusioni relative alle prestazioni vengono formulate solo dopo l'analisi
dei dati raccolti; le campagne definiscono quindi il protocollo di confronto,
non il risultato atteso.


Una sintesi delle campagne è riportata in
\cref{tab:campagne-sperimentali}.

\begin{table}[H]
\centering
\footnotesize
\caption{Sintesi delle campagne sperimentali.}
\label{tab:campagne-sperimentali}
\begin{tabular}{@{}clp{0.54\linewidth}@{}}
\toprule
\textbf{ID} & \textbf{Variabile studiata} & \textbf{Obiettivo} \\
\midrule

C1 &
$M$, $N$, rapporto $M{:}N$ &
Valutare l'effetto della dimensione del problema e della forma della matrice
su configurazioni CPU e GPU. \\

C2 &
$k$ e backend &
Confrontare le diverse implementazioni al variare del numero di colonne del
multivettore e verificare il comportamento del percorso generico. \\

C3 &
specializzazione di $k$ &
Isolare il beneficio della specializzazione a compile-time confrontando,
sullo stesso problema, versione specializzata e versione generica del kernel
CPU. \\

C4 &
numero di processi $P$ &
Studiare lo strong scaling mantenendo costante la dimensione globale del
problema. \\

C5 &
numero di processi $P$ e dimensione globale &
Studiare il weak scaling mantenendo costante la dimensione del blocco locale
assegnato a ciascun processo. \\

C6 &
forma $P_r \times P_c$ &
Valutare l'effetto della forma della griglia MPI mantenendo costante il
numero complessivo di processi. \\

C7 &
precisione FP32/FP64 &
Confrontare le due precisioni numeriche e correlare le prestazioni alle
risorse utilizzate dai kernel CUDA. \\

\bottomrule
\end{tabular}
\end{table}